\chapter{Resumen de Postulados y Teoremas (Notación Ingenieril)}
En esta sección se recopilan los postulados de Huntington y los teoremas principales derivados, transcritos a la notación propia del álgebra de conmutación y la lógica digital ($+,$ $\cdot,$ $0,$ $1,$ $\overline{a}$), concebidos como hoja de referencia rápida.

\subsection*{Pre-Axiomas de la Estructura}

\setcounter{tcb@cnt@preaxioma}{0}
\begin{quote}\textbf{Preaxioma (Estructura de Conjunto):}\label{esconj_eng}
Se requiere que se defina sobre un conjunto (por ejemplo, $\mathbb{B}_2 = \{0, 1\}$).
\end{quote}

\setcounter{tcb@cnt@preaxioma}{1}
\begin{quote}\textbf{Preaxioma (Constantes Lógicas):}\label{constantes_eng}
Este conjunto contiene dos constantes fundamentales: $0$ (falso) y $1$ (verdadero).
\end{quote}

\setcounter{tcb@cnt@preaxioma}{2}
\begin{quote}\textbf{Preaxioma (Operaciones Binarias Internas):}\label{opbinint_eng}
Se definen dos operaciones binarias internas, la suma ($+$) y el producto ($\cdot$):
\begin{align*}
+ &: \mathbb{B}_2 \times \mathbb{B}_2 \to \mathbb{B}_2 \\
\cdot &: \mathbb{B}_2 \times \mathbb{B}_2 \to \mathbb{B}_2
\end{align*}
\end{quote}

\setcounter{tcb@cnt@preaxioma}{4}
\begin{quote}\textbf{Preaxioma (Existencia y Unicidad de Imagen):}\label{exist_unic_eng}
Para cada par de elementos del conjunto, las operaciones $+$ y $\cdot$ siempre producen un resultado que también pertenece al conjunto, y ese resultado es siempre único.
\end{quote}

\subsection*{Postulados de Huntington}

\setcounter{tcb@cnt@postulado}{0}
\begin{quote}\textbf{Postulado (Elemento neutro):}\label{neutro_eng}
\[ a + 0 = a \qquad \text{y} \qquad a \cdot 1 = a \]
\end{quote}

\setcounter{tcb@cnt@postulado}{2}
\begin{quote}\textbf{Postulado (Conmutatividad):}\label{conmut_eng}
\[ a + b = b + a \qquad \text{y} \qquad a \cdot b = b \cdot a \]
\end{quote}

\setcounter{tcb@cnt@postulado}{4}
\begin{quote}\textbf{Postulado (Distributividad):}\label{distrib_eng}
\[ a \cdot (b + c) = (a \cdot b) + (a \cdot c) \qquad \text{y} \qquad a + (b \cdot c) = (a + b) \cdot (a + c) \]
\end{quote}

\setcounter{tcb@cnt@postulado}{6}
\begin{quote}\textbf{Postulado (Complementario):}\label{comp_eng}
Para cada elemento $a,$ existe un complemento $\overline{a}$ tal que:
\[ a + \overline{a} = 1 \qquad \text{y} \qquad a \cdot \overline{a} = 0 \]
\end{quote}

\subsection*{Teoremas Fundamentales}

\setcounter{tcb@cnt@teorema}{0}
\begin{quote}\textbf{Teorema (Unicidad de los elementos neutros):}\label{unicidad_neutros_eng}
El elemento neutro para la suma ($0$) y para el producto ($1$) son únicos.
\end{quote}

\setcounter{tcb@cnt@teorema}{1}
\begin{quote}\textbf{Teorema (Idempotencia):}\label{idempotencia_eng}
\[ a + a = a \qquad \text{y} \qquad a \cdot a = a \]
\end{quote}

\setcounter{tcb@cnt@teorema}{2}
\begin{quote}\textbf{Teorema (Elementos absorbentes):}\label{absorbentes_eng}
\[ a + 1 = 1 \qquad \text{y} \qquad a \cdot 0 = 0 \]
\end{quote}

\setcounter{tcb@cnt@teorema}{5}
\begin{quote}\textbf{Teorema (Propiedades de absorción):}\label{absorcion_eng}
\[ a + (a \cdot b) = a \qquad \text{y} \qquad a \cdot (a + b) = a \]
\end{quote}

\setcounter{tcb@cnt@teorema}{11}
\begin{quote}\textbf{Teorema (Leyes de De Morgan):}\label{morgan_eng}
\[ \overline{a + b} = \overline{a} \cdot \overline{b} \qquad \text{y} \qquad \overline{a \cdot b} = \overline{a} + \overline{b} \]
\end{quote}

\setcounter{tcb@cnt@teorema}{10}
\begin{quote}\textbf{Teorema (Involución (Doble negación)):}\label{involucion_eng}
\[ \overline{\overline{a}} = a \]
\end{quote}

\setcounter{tcb@cnt@teorema}{12}
\begin{quote}\textbf{Teorema (Asociatividad):}\label{asociatividad_eng}
\[ a + (b + c) = (a + b) + c \qquad \text{y} \qquad a \cdot (b \cdot c) = (a \cdot b) \cdot c \]
\end{quote}

\setcounter{tcb@cnt@teorema}{9}
\begin{quote}\textbf{Teorema (Unicidad del complemento):}\label{unic_comp_eng}
El complemento $\overline{a}$ de un elemento $a$ es único.
\end{quote}

\setcounter{tcb@cnt@teorema}{6}
\begin{quote}\textbf{Teorema (Otras propiedades equivalentes):}\label{otras_prop_eng}
\begin{itemize}
    \item \textbf{Orden de retículo:} $a + b = a \iff a \cdot b = b$
    \item \textbf{Equivalencia de operaciones:} $a + b = a \cdot b \implies a = b$
    \item \textbf{Cancelación:} $(a + b = a + c \text{ y } a \cdot b = a \cdot c) \implies b = c$
\end{itemize}
\end{quote}

\setcounter{tcb@cnt@definicion}{0}
\begin{quote}\textbf{Definicion (Generalización a $n$ variables):}\label{gen_n_vars_eng}
Las operaciones disyunción y conjunción pueden extenderse a un número $n$ de variables mediante los símbolos sumatorio y productorio:
\[ \sum_{i=1}^{n} x_i = x_1 + x_2 + \dots + x_n \]
\[ \prod_{i=1}^{n} x_i = x_1 \cdot x_2 \cdot \dots \cdot x_n \]
\end{quote}

\setcounter{tcb@cnt@teorema}{3}
\begin{quote}\textbf{Teorema (Casos de Álgebra Trivial):}\label{trivial_eng}
Si $0 = 1,$ o si existe algún elemento tal que $\overline{a} = a,$ entonces el álgebra contiene un único elemento (álgebra trivial).
\end{quote}


\setcounter{tcb@cnt@teorema}{13}
\begin{quote}\textbf{Teorema (Teorema de Adyacencia (Expansión de Shannon)):}\label{adyacencia_eng}
\[ a = (a \cdot b) + (a \cdot \overline{b}) \qquad \text{y} \qquad a = (a + b) \cdot (a + \overline{b}) \]
\end{quote}

\setcounter{tcb@cnt@teorema}{14}
\begin{quote}\textbf{Teorema (Teorema de Reducción (Absorción Fuerte)):}\label{reduccion_eng}
\[ a + (\overline{a} \cdot b) = a + b \qquad \text{y} \qquad a \cdot (\overline{a} + b) = a \cdot b \]
\end{quote}

\setcounter{tcb@cnt@teorema}{15}
\begin{quote}\textbf{Teorema (Teorema del Consenso (Quine)):}\label{consenso_eng}
\[ (a \cdot b) + (\overline{a} \cdot c) + (b \cdot c) = (a \cdot b) + (\overline{a} \cdot c) \]
\[ (a + b) \cdot (\overline{a} + c) \cdot (b + c) = (a + b) \cdot (\overline{a} + c) \]
\end{quote}

\subsection*{Comportamiento de Operadores Derivados}

\setcounter{tcb@cnt@teorema}{13}
\begin{quote}\textbf{Teorema (Idempotencia cruzada (NAND/NOR)):}\label{idemp_cruzada_eng}
\[ a \uparrow a = \overline{a} \qquad \text{y} \qquad a \downarrow a = \overline{a} \]
\end{quote}

\setcounter{tcb@cnt@teorema}{14}
\begin{quote}\textbf{Teorema (Generación de AND y OR):}\label{gen_inf_sup_eng}
\[ a \cdot b = \overline{a \uparrow b} = (a \uparrow b) \uparrow (a \uparrow b) \qquad \text{y} \qquad a + b = \overline{a \downarrow b} = (a \downarrow b) \downarrow (a \downarrow b) \]
\end{quote}

\setcounter{tcb@cnt@teorema}{15}
\begin{quote}\textbf{Teorema (Generación cruzada):}\label{gen_cruzada_eng}
\[ a + b = \overline{a} \uparrow \overline{b} = (a \uparrow a) \uparrow (b \uparrow b) \qquad \text{y} \qquad a \cdot b = \overline{a} \downarrow \overline{b} = (a \downarrow a) \downarrow (b \downarrow b) \]
\end{quote}

\setcounter{tcb@cnt@teorema}{16}
\begin{quote}\textbf{Teorema (Conmutatividad):}\label{conmut_deriv_eng}
\[ a \uparrow b = b \uparrow a \qquad \text{y} \qquad a \downarrow b = b \downarrow a \]
\end{quote}

\setcounter{tcb@cnt@teorema}{18}
\begin{quote}\textbf{Teorema (Comportamiento con las constantes):}\label{constantes_deriv_eng}
\begin{align*}
    a \uparrow 1 &= \overline{a} \qquad & a \downarrow 0 &= \overline{a} \\
    a \uparrow 0 &= 1 \qquad & a \downarrow 1 &= 0
\end{align*}
\end{quote}

\setcounter{tcb@cnt@teorema}{19}
\begin{quote}\textbf{Teorema (Ausencia de Asociatividad):}\label{no_asoc_deriv_eng}
\[ (a \uparrow b) \uparrow c \neq a \uparrow (b \uparrow c) \qquad \text{y} \qquad (a \downarrow b) \downarrow c \neq a \downarrow (b \downarrow c) \]
\end{quote}

\setcounter{tcb@cnt@definicion}{6}
\begin{quote}\textbf{Definicion (NAND y NOR de 3 entradas):}\label{n_entradas_eng}
\[ \uparrow(a,b,c) = \overline{a \cdot b \cdot c} \qquad \text{y} \qquad \downarrow(a,b,c) = \overline{a + b + c} \]
\end{quote}

\setcounter{tcb@cnt@teorema}{20}
\begin{quote}\textbf{Teorema (NAND/NOR múltiple vs cascada binaria):}\label{multiple_vs_binaria_eng}
\begin{align*}
    \uparrow(a,b,c) &\neq (a \uparrow b) \uparrow c \qquad & \uparrow(a,b,c) &\neq a \uparrow (b \uparrow c) \\
    \downarrow(a,b,c) &\neq (a \downarrow b) \downarrow c \qquad & \downarrow(a,b,c) &\neq a \downarrow (b \downarrow c)
\end{align*}
\end{quote}
\subsection*{Comportamiento de los Operadores XOR y XNOR}

\setcounter{tcb@cnt@definicion}{4}
\begin{quote}\textbf{Definicion (Definición de XOR y XNOR):}\label{def_xor_xnor_eng}
\[ a \oplus b = (a \cdot \overline{b}) + (\overline{a} \cdot b) \qquad \text{y} \qquad a \odot b = \overline{a \oplus b} = (a \cdot b) + (\overline{a} \cdot \overline{b}) \]
\end{quote}

\setcounter{tcb@cnt@teorema}{22}
\begin{quote}\textbf{Teorema (Conmutatividad):}\label{conmut_xor_eng}
\[ a \oplus b = b \oplus a \qquad \text{y} \qquad a \odot b = b \odot a \]
\end{quote}

\setcounter{tcb@cnt@teorema}{23}
\begin{quote}\textbf{Teorema (Elementos Neutros e Inversores):}\label{neutros_xor_eng}
\begin{align*}
    a \oplus 0 &= a \qquad & a \odot 1 &= a \\
    a \oplus 1 &= \overline{a} \qquad & a \odot 0 &= \overline{a}
\end{align*}
\end{quote}

\setcounter{tcb@cnt@teorema}{24}
\begin{quote}\textbf{Teorema (Elemento Inverso de sí mismo (Grupo Abeliano)):}\label{idemp_nula_eng}
\[ a \oplus a = 0 \qquad \text{y} \qquad a \odot a = 1 \]
\end{quote}

\setcounter{tcb@cnt@teorema}{25}
\begin{quote}\textbf{Teorema (Propiedades de Negación):}\label{neg_xor_eng}
\[ \overline{a \oplus b} = \overline{a} \oplus b = a \oplus \overline{b} = a \odot b \]
\[ \overline{a \odot b} = \overline{a} \odot b = a \odot \overline{b} = a \oplus b \]
\end{quote}

\setcounter{tcb@cnt@teorema}{26}
\begin{quote}\textbf{Teorema (Asociatividad y Generalización):}\label{asoc_gen_xor_eng}
Ambos operadores son asociativos:
\[ a \oplus (b \oplus c) = (a \oplus b) \oplus c \qquad \text{y} \qquad a \odot (b \odot c) = (a \odot b) \odot c \]
Lo cual permite su generalización a un número arbitrario $n$ de entradas:
\[ \bigoplus_{i=1}^{n} x_i = x_1 \oplus x_2 \oplus \dots \oplus x_n \qquad \text{y} \qquad \bigodot_{i=1}^{n} x_i = x_1 \odot x_2 \odot \dots \odot x_n \]
\end{quote}

\setcounter{tcb@cnt@teorema}{28}
\begin{quote}\textbf{Teorema (Distributividad con el producto y la suma):}\label{dist_xor_eng}
\[ a \cdot (b \oplus c) = (a \cdot b) \oplus (a \cdot c) \qquad \text{y} \qquad a + (b \odot c) = (a + b) \odot (a + c) \]
\end{quote}

\subsection*{Tablas de Verdad Bivaluadas}

Resumen de las tablas de operación del álgebra de Boole para el caso de dos elementos ($B=\{0,1\}$), transcritas al lenguaje ingenieril:

\begin{table}[h!]
\centering
\renewcommand{\arraystretch}{1.2}
\begin{tabular}{c||c}
\hline
$a$ & $\overline{a}$ \\
\hline\hline
$0$ & $1$ \\
$1$ & $0$ \\
\hline
\end{tabular}
\qquad
\begin{tabular}{cc||c|c|c|c|c|c}
\hline
$a$ & $b$ & $a + b$ & $a \cdot b$ & $a \uparrow b$ & $a \downarrow b$ & $a \oplus b$ & $a \odot b$ \\
\hline\hline
$0$ & $0$ & $0$ & $0$ & $1$ & $1$ & $0$ & $1$ \\
$0$ & $1$ & $1$ & $0$ & $1$ & $0$ & $1$ & $0$ \\
$1$ & $0$ & $1$ & $0$ & $1$ & $0$ & $1$ & $0$ \\
$1$ & $1$ & $1$ & $1$ & $0$ & $0$ & $0$ & $1$ \\
\hline
\end{tabular}
\end{table}

\makeatother
